\chapter{One-shot luck and coin-luck}\label{ch:coin}

\textit{Ingredients from Luck:} the definition $L(x) = |\Omega(x)| +
\tfrac{1}{2}|\omega(x)|$; the empirical estimator $\ell$ and its variance
bound $L(1-L)/n \le 1/(4n)$; numerical luck via simulation.

\section{The 36 heads problem}

Someone tosses a fair coin 36 times in a row and reports 36 heads.  Should
you believe the coin is fair?

Under the fair-coin model, the probability of \emph{any} specific 36-flip
sequence is $2^{-36} \approx 1.5 \times 10^{-11}$.  All sequences are
equally likely.  Naive ``rare $\Rightarrow$ surprising'' reasoning
collapses: every sequence is rare.  What separates 36-heads from
HTHHTH\ldots is not its probability but its \emph{rank} in the
distribution of $-\log p$ values --- and, equivalently, its luck.

For 36 flips of a fair coin, there are $2^{36}$ outcomes, all of equal
probability.  $\Omega(x) = \emptyset$ (no outcome is more probable than
$x$), $\omega(x)$ is the entire space, and $L(x) = 1/2$.  Every outcome
has the same luck, and luck does not flag 36-heads as anomalous --- because
under the fair-coin model, it isn't.

The interesting computation is under the \emph{alternative}: a biased
coin.  If $p = 0.6$, the probability of 36 heads is $0.6^{36} \approx
1.0 \times 10^{-8}$, but the probability of getting 36 heads is now the
maximum-probability outcome, so $L(\text{36 H}) = 1/2 \cdot p(36\text{H})
\approx 5 \times 10^{-9}$.  Under $p=0.6$, 36 heads is the
\emph{luckiest} possible result.

The point: luck is always relative to a model.  ``36 heads is incredible''
is a statement about which model you were holding.

\section{Coin-luck as a unit}

The natural way to talk about luck is in coin flips.  Define
\[
  C(x) = \log_2 \frac{L(x)}{1 - L(x)}\,,
\]
the \emph{coin-luck} of $x$.  $C = 0$ is typical; $C = +n$ means $x$ is in
the rare tail by roughly $n$ coin-flip surprise units; $C = -n$ means $x$
is so common that finding it is no surprise at all.

\begin{example}{36 heads, biased coin.}
For $p=0.6$, $L \approx 5 \times 10^{-9}$, so $C \approx +27.5$.  Reading
this aloud: ``36 heads on a $p=0.6$ coin is about 27 coin-flips' worth of
surprise.''  The reader does not need a probability scale; the coin-flip
unit is its own scale.
\end{example}

The coin-luck framing also explains why people are bad at intuition for
small probabilities.  Probabilities below $10^{-3}$ all feel the same
(``small''); coin-lucks above $10$ are still distinguishable from coin-lucks
above $20$.  The logarithm is doing the work the brain refuses to.

\section{When to use $\ell$, $L$, and $C$ in conversation}

\begin{itemize}
\item $\ell$: when reporting a number computed from data, with a noise
  bar.
\item $L$: when speaking analytically, in proofs and thresholds.
\item $C$: when explaining to a non-statistician, or comparing rare events
  across orders of magnitude.
\end{itemize}

\section{The Read lever}

Every \emph{Read} chapter in this book applies the same recipe to a
specific domain:
\begin{enumerate}
\item Identify the model under which luck is being read.
\item Compute $\ell$ (or $L$, when tractable) for the observation.
\item Convert to coin-luck $C$ for interpretation.
\item Decide: act, investigate, or update the model.
\end{enumerate}

The decision step is what changes by chapter --- a clinician acts on
$|C| > 4$ differently than a quality-control engineer or a fraud analyst.
The first three steps are universal.

\section{Worked example: detecting a loaded die}
\todo{Six-sided die rolled $n$ times; observed face counts compared to
fair-die luck; coin-luck threshold for declaring loaded.  Use this to
show: (a) what $\ell$ looks like for a discrete model, (b) how the noise
bar $\sqrt{L(1-L)/n}$ translates into a coin-luck uncertainty, (c) when
the $|\Omega| + \tfrac{1}{2}|\omega|$ tie correction matters.}

\section{Common mistakes}

\paragraph{Confusing $L$ with $p$.}  An outcome with $p = 10^{-9}$ is not
necessarily lucky; if every outcome has $p = 10^{-9}$, $L = 1/2$.  Luck is
\emph{rank-based on probability}, not probability itself.

\paragraph{Reading luck without naming the model.}  ``This is one in a
million'' is meaningless without specifying \emph{under what
distribution}.  Every coin-luck statement has a hidden ``\dots under model
$M$'' clause.

\paragraph{Treating $C$ as a $z$-score.}  Coin-luck is logistic, not
Gaussian.  $|C| > 3$ is not a $0.27\%$ tail; it is closer to $11\%$ for a
single observation under a correctly-fit model.  The chapters in
Part~II calibrate the right thresholds for each application.
